\section{Discussion: What Generalises}
\label{sec:lessons}

Nothing in this project required a novel technique. What it required, repeatedly,
was noticing that a question had not been asked. This section extracts the parts
that transfer.

\subsection{Search order dominates search effort}

The defect was in a component that was never varied, and it was found within one
attempt of finally varying it. The preceding weeks were spent varying components
that could not have been responsible.

The failure of process is identifiable in advance. \textbf{A component chosen at
the start of a project because it worked becomes invisible to the project.} It is
not on the list of suspects because it is part of the definition of the setup.

The practical rule follows from delta debugging~\cite{zeller2002}, applied one
level up from where it is usually applied:

\begin{quote}
\textbf{Before bisecting within a component, enumerate the components you have
never bisected across.} Any element of the stack that has held one value for the
whole investigation belongs on the suspect list \emph{ahead of} the elements you
have already varied without success, not behind them.
\end{quote}

A corroborating signal was available and misread: eleven configuration changes
inside the translation layer produced no effect at all (Table~\ref{tab:shaderelim}).
An exhausted configuration surface is evidence that the component is the wrong
one to configure---not an invitation to configure it harder.

\subsection{Differential diagnosis beats hypothesis generation}

Twice in this project, two failures shared a symptom and were separated by one
cheap observation:

\begin{itemize}
\item Extraction: process alive, nothing happening. \emph{Is a core busy, and did
      it write any bytes?} Two seconds of observation, two entirely different
      root causes (Table~\ref{tab:diff}).
\item Rendering: window open, nothing happening. \emph{Ask the window manager
      what windows exist.} One command, and an invisible modal dialog becomes a
      driver crash address.
\end{itemize}

In both cases the correct move was to \emph{enumerate the state} rather than to
generate hypotheses about it. Hypothesis generation is unbounded; state
enumeration is finite and usually one command.

\subsection{Prefer pipelines whose stages publish their preconditions}

The reassembly (Section~\ref{sec:reassembly}) succeeded on the first attempt
across ten stages and 54\,GB because every stage could be checked against a
manifest read \emph{before} it ran. Five checks, five exact matches, and any
error would have surfaced at its own stage.

Where a format offers this, it is the cheapest correctness guarantee available:
a checksum you did not have to compute and did not have to trust yourself to
choose.

\subsection{Instruments lie, and only negative testing catches it}

Every measurement tool built in this project produced at least one confident
false reading before being corrected:

\begin{itemize}
\item The assertion counter reported zero because \verb|grep -c| both prints zero
      and fails.
\item The override verifier reported success on runs that overrode nothing, for
      the same reason.
\item The process matcher reported the game dead because \verb|comm| truncates at
      fifteen characters.
\item The graphics layer reported an AMD adapter while running on NVIDIA, by
      design.
\item The launcher's own notification claimed NVIDIA while using the integrated
      GPU.
\item The recorded conclusion ``NVIDIA Vulkan is non-functional'' was an artefact
      of a system-wide policy the test was unaware of.
\end{itemize}

\textbf{The discipline that catches all six is the same: run every instrument
against an input whose answer you already know, and specifically against a
\emph{failing} input.} A health check that has only ever been shown reporting
health has not been tested.

\subsection{``Impossible'' has exactly one legitimate use}

This paper contains one genuine impossibility claim---the hypervisor-based
countermeasure of Section~\ref{sec:contrast}, which fails for two independent
structural reasons---and it is the only place the word is used.

Elsewhere, several things that were recorded as impossible turned out to be
merely unexamined: a compiled installer script that ``cannot be read'' was read
by parsing the format; a driver that was ``dead'' was disabled by an environment
variable; an archive that used ``an unsupported codec'' was missing a
configuration file.

\begin{quote}
\textbf{Separate three claims that are easily conflated: this cannot exist; this
would be difficult; I do not currently know how.} Only the first is a stopping
condition, and it requires a structural argument, not a failed attempt.
\end{quote}

\subsection{The strong form of a negative result}

Section~\ref{sec:denuvo} makes a claim of absence---that no commercial anti-tamper
system was ever in this build---and supports it with a cryptographic identity
rather than an accumulation of heuristics.

This is worth generalising because negative results about binaries are usually
argued the weak way: entropy, section names, string searches, absence of known
signatures. Each of those is defeasible individually, and stacking defeasible
arguments does not produce an indefeasible one. A validated signature over the
current bytes does, because it converts ``we looked and found nothing'' into
``the file is provably identical to the one the publisher signed, so there is
nothing to find.''

The corollary is a cost argument: \textbf{the strong test was also the cheapest
one.} It was performed after the weak ones, which is the wrong order.

\subsection{Threats to validity}

\begin{itemize}
\item \textbf{Single system.} Every result comes from one laptop with one hybrid
      graphics configuration. The reassembly and forensic findings should
      transfer directly; the graphics findings are specific to this driver
      version, this translation layer, and this vendor pair.
\item \textbf{The root cause is localised, not explained.} We show that one build
      of the translation layer produces shader code that crashes a specific
      driver's compiler and that another does not. We did not identify the
      offending construct, and did not reduce it to a minimal reproducer. That
      would be the natural next step and would produce a filable driver bug.
\item \textbf{The unresolved contradiction of Section~\ref{sec:shaders}} is
      genuinely unresolved. The partial improvement from shader substitution has
      no confirmed mechanism.
\item \textbf{Timing figures} are wall-clock on a loaded desktop system, not a
      controlled benchmark.
\end{itemize}

\subsection{Related work}

Entropy analysis for packing detection follows Lyda and
Hamrock~\cite{lyda2007}; the windowing refinement and the known false-positive
behaviour are discussed at length in the packer-detection
literature~\cite{sokpacker}. Authenticode structure and validation semantics
are specified by Microsoft~\cite{authenticode} over PKCS\#7/CMS~\cite{rfc5652}
and the PE format~\cite{pecoff}. The bisection discipline in
Section~\ref{sec:lessons} is delta debugging~\cite{zeller2002}, and the general
approach to systematic debugging follows Zeller~\cite{zellerbook}. The
compatibility stack under test is Wine~\cite{wine} with Proton~\cite{proton},
vkd3d-proton~\cite{vkd3d}, and DXVK~\cite{dxvk}, targeting
Vulkan~\cite{vulkan} and SPIR-V~\cite{spirv}. The reassembly toolchain uses
LZ4~\cite{lz4}, VCDIFF~\cite{rfc3284}, HDiffPatch~\cite{hdiffpatch}, and Ogg
Vorbis~\cite{vorbis}; the installer format is Inno Setup~\cite{innosetup} and
the object-namespace behaviour is documented in Windows
Internals~\cite{ntobj}. Vendor dispatch is handled by libglvnd~\cite{glvnd}.

To our knowledge the specific combination reported here---a signature-based
negative determination of anti-tamper presence, a manual reconstruction of a
delta-compressed redistribution without its installer, and a cross-vendor shader
divergence causing a driver-side compiler fault under translation---has not been
described together in one place.
